知性は長らく、能力・性質・蓄積された知識・正しさの所有 ── いずれも何らかの \textbf{停止構造} ── として論じられてきた。本論文はこの枠組みを退け、知性を「ある構造的条件のもとでのみ持続する力学的現象」として定式化する。

具体的には三つの最小公理 ── 連結された状態空間、時間的順序付け、グローバル制約 $\Phi$ ── を提示し、これらが知覚・知識・知性という三つの層を区別するための最小条件として働くことを示す。さらに粘菌(\emph{Physarum polycephalum})の網状動態と Transformer の層遷移という substrate の根本的に異なる二系を、これら三公理のもとに構造的に対応するものとして記述する。とりわけ Tero, Kobayashi \& Nakagaki (2007) による Physarum solver の flow-reinforcement 動力学が、本論文における $\Phi$ の創発性(局所更新規則への非還元性)と同型の構造に、独立な実証研究路線から到達していたことを示す。

この構造から、Gettier 問題は知識の定義の欠陥ではなく、「停止構造に動的性質を担保させようとする層間ミスマッチ」として読み直される。さらに、正しさの累積はそれ自体が分岐数の爆発を生み、個体内だけでは知性が安定化しない ── すなわち \textbf{停止演算子は個体内には存在しない} ── ことが帰結する。

この帰結が要請する不可避の越境点を、本論文では \emph{the differential horizon of intelligence} と命名する。これは VED Vol.4 における \DH{} の別個の複製ではなく、同一の記述境界が知性層の観測窓に局所的に現れる horizon-form である。物理層・生命層(Bio-IFGT)・意識層(意識編)・知性層の四層対応を提示し、本論文の射程の外で扱われる後続課題への接続点を明示する。

本論文の理論的基盤は Vortical Enclosure Dynamics (VED) および Information Field Geometry Theory (IFGT) であり、それぞれが提供する記述語彙のもとで構造論的議論を展開する。
